\section{Conclusion}
\label{sec:conclusion}

This report combines a historical-data audit, a predictive LOS interface,
five metaheuristics and a room-level simulation-optimisation framework.
It formalises separate session-and-bed and operational room objectives,
describes the shared decoder and acceptance loop, and retains the provenance
of each experimental contribution.

The LOS results show that session granularity can change algorithm rankings
at unchanged total capacity. The hybrid attains the lowest mean cost among
four split-session configurations, with a 5.70\% gap to an overload lower
bound under unequal search budgets. Upper-endpoint planning does not
consistently reduce scenario-mean cost. In the room campaign, all searches
reach the small fixture's exact initial allocation optimum, while larger
predicted improvements need not survive execution. Reactive scheduling
improves completions in 16 of 144 closure pairs, ties in 83 and worsens them
in 45. The separate oracle experiment illustrates that increased throughput
can require increased overtime even with perfect duration information.
These are descriptive findings under explicit modelling assumptions, not
a universal method ranking or evidence of clinical effectiveness.

Further work should reconcile LOS and bed-accounting conventions, evaluate
current learned room-duration models on held-out workloads, compare total
computation as well as per-request budgets, and vary loads, disruptions and
sequencing policies. Integrating ward dynamics, staff and clinical constraints
would connect the current experiments more closely to admission-date decision
support. Such extensions require additional data and validation; they are
not implied by the successful software integration.
